\documentclass{article}
\usepackage{graphicx}
\usepackage{amsfonts}
\usepackage{amsmath}
\usepackage{amsthm}
\usepackage{tikz}

\theoremstyle{definition}
\newtheorem{theorem}{Theorem}[section]
\newtheorem{corollary}[theorem]{Corollary}
\newtheorem{proposition}[theorem]{Proposition}
\newtheorem{lemma}[theorem]{Lemma}
\newtheorem{definition}[theorem]{Definition}

\theoremstyle{remark}
\newtheorem{remark}[theorem]{Remark}
\newtheorem{example}[theorem]{Example}
\newtheorem{claim}[theorem]{Claim}

\title{Complex Vector Space}
\author{Sarah Liu}
\date{January 2026}

\begin{document}

\maketitle

\section{Complex Number}

\begin{definition}[Complex Numbers]
    The set of complex number
    \[
    \mathbb{C}:=\{a+bi:a,b\in\mathbb{R}\}
    \]
    where $i^2=-1$.
\end{definition}

\begin{definition}
    The absolute value of $z$ is 
    \[
    |z|=\sqrt{a^2+b^2}=\sqrt{z\cdot\bar{z}}=|\bar{z}|
    \]
    and $\theta$ is the argument. We could write
    \[
    z=|z|(\cos\theta+i\sin\theta)
    \]
\end{definition}

\begin{theorem}[Moivie's Formula]
    \[
    z^n=|z|^n(\cos n\theta+i\sin n\theta)
    \]
\end{theorem}

\section{Complex Polynomials}
\begin{theorem}
    Let $p(z)=a_n z^n+\cdots+a_1 z+a_0$. Then $p(z)$ has $n$ roots counted with multiplicities.
\end{theorem}

\begin{proposition}
    Let $p(z)=a_n z^n+\cdots+a_1 z+a_0$, where $a_0,a_1,\dots,a_n\in\mathbb{R}$.
    If $p(z)=0$, then $p(\bar{z})=0$.
\end{proposition}

\section{Complex Vector Space}
\begin{definition}
    A field is a set $F$ with two operation $(+,\cdot)$ satisfying that for $x,y,z\in F$
    \begin{enumerate}
        \item $x+y=y+x$
        \item $(x+y)+z=x+(y+z)$
        \item There exists $0\in F$, $0+x=x$
        \item There exists $-x\in F$, $x+(-x)=0$
        \item $x\cdot y=y\cdot x$
        \item $(x\cdot y)\cdot z=x\cdot (y\cdot z)$
        \item There exists $1\in F$, $x\cdot 1=x$
        \item If $x\neq 0$, then there exists $x^{-1}\in F$, $x\cdot x^{-1}=1$
        \item $(x+y)\cdot z=x\cdot z+y\cdot z$
    \end{enumerate}
\end{definition}

\begin{definition}
    A Hertian inner product on $V$ is $\langle,\rangle:V\times V\rightarrow \mathbb{C}$ satisfying for all $\vec{u},\vec{v},\vec{w}\in V$
    \begin{enumerate}
        \item $\langle a\vec{u}+b\vec{v},\vec{w}\rangle=a\langle\vec{u},\vec{w}\rangle+b\langle\vec{v},\vec{w}\rangle$
        \item $\langle u,v\rangle=\overline{\langle v,u\rangle}$
        \item $\langle \vec{v},\vec{v}\rangle\geq 0$
    \end{enumerate}
\end{definition}

\begin{remark}
    The second entry 
    \[
    \langle \vec{u},a\vec{v}\rangle =\overline{\langle a\vec{v},\vec{u}\rangle}=\bar{a}\overline{\langle\vec{v},\vec{u}\rangle}=\bar{a}\langle\vec{u},\vec{v}\rangle
    \]
\end{remark}

\begin{definition}[Standard Hermition Product]
\[
\left\langle\begin{bmatrix}
z_1 \\
z_2 \\
\vdots \\
z_n
\end{bmatrix},
\begin{bmatrix}
w_1 \\
w_2 \\
\vdots \\
w_n
\end{bmatrix}
\right\rangle=z_1\bar{w_1}+z_2\bar{w_2}+\cdots+z_n\bar{w_n}
\]
    
\end{definition}

\section{Self-Adjoint}
\begin{definition}
    Let $A$ be an $n\times n$ matrix, the adjoint of $A$ is $A^*=(\overline{A})^T$.
    $A$ is self-adjoint if $A=A^*$.
\end{definition}

\begin{definition}
    Let $T:\mathbb{C}\rightarrow\mathbb{C}$ be a linear transformation.
    Let $\alpha$ be an orthonormal basis of $\mathbb{C}$. 
    The adjoint of linear transformation $T^*$ is described by the matrix
    \[
    [T^*]^{\alpha}_{\alpha}=(\overline{[T]^{\alpha}_{\alpha}})^T
    \]
\end{definition}

\begin{theorem}
    For all $\vec{z},\vec{w}\in\mathbb{C}^n$, $\langle T(\vec{z}),\vec{w}\rangle=\langle \vec{z},T^*(\vec{w})\rangle$
\end{theorem}

\begin{proof}
    
\end{proof}

\begin{theorem}
    Let $T:\mathbb{C}^n\rightarrow\mathbb{C}^n$ be a self-adjoint linear transformation. 
    If $\lambda$ be an eigenvalue for $T$, then $\lambda\in\mathbb{R}$.
\end{theorem}

\begin{proof}
    Let $\vec{v}$ be an eigenvalue satisfying $\langle\vec{v},\vec{v}\rangle=1$.
    \[
    \lambda=\lambda\langle\vec{v},\vec{v}\rangle=\langle\lambda\vec{v},\vec{v}\rangle=\langle T(\vec{v}),\vec{v}\rangle=\langle\vec{v},T(\vec{v})\rangle=\langle\vec{v},\lambda\vec{v}\rangle=\bar{\lambda}\langle\vec{v},\vec{v}\rangle=\bar{\lambda}
    \]
\end{proof}

\begin{proposition}
    If $\vec{u},\vec{v}$ are eigenvectors of self-adjoint linear transformation $T$ with eigenvalues $\mu\neq\lambda$, then $\langle\vec{u},\vec{v}\rangle=0$.
\end{proposition}

\begin{proof}
    \[
    \mu \langle\vec{u},\vec{v}\rangle=\langle\mu\vec{u},\vec{v}\rangle=\langle T(\vec{u}),\vec{v}\rangle=\langle\vec{u},T(\vec{v})\rangle=\langle\vec{u},\lambda \vec{v}\rangle=\lambda\langle\vec{u},\vec{v}\rangle
    \]
    \[
    (\mu-\lambda)\langle\vec{u},\vec{v}\rangle=0
    \]
    \[
    \langle\vec{u},\vec{v}\rangle=0
    \]
\end{proof}

\begin{theorem}
    Let $T:\mathbb{C}^n\rightarrow\mathbb{C}^n$ be a self-adjoint linear transformation. There exists an orthonormal basis consisting of eigenvectors if and only if it is diagonalizable.
\end{theorem}

\end{document}